\documentclass[12pt,a4paper]{article}
%\usepackage[utf-8]{inputenc}
\usepackage[margin=1in]{geometry}
\usepackage{amsmath}
\usepackage{amssymb}
\usepackage{graphicx}
\usepackage{array}
\usepackage{booktabs}
\usepackage{xcolor}
\usepackage{fancyhdr}
\usepackage{hyperref}

\pagestyle{fancy}
\fancyhf{}
\rhead{GED Math Solutions}
\lhead{Graphs and Averages}
\cfoot{\thepage}

\title{\textbf{GED Mathematics Solutions}\\[0.3cm]\Large{Reading Graphs, Tables, and Computing Averages}\\[0.2cm]\normalsize{Complete Worked Solutions for 50 Practice Questions}}
\author{}
\date{May 2026}

\begin{document}

\maketitle

\tableofcontents

\newpage

\section{Part A: Reading and Interpreting Tables (Solutions 1-10)}

\subsection*{Solution 1}
\textbf{Question:} The table shows quarterly sales (in thousands) for a retail store in 2025. What was the total sales for the year?

\textbf{Given:} Q1 = \$45,000, Q2 = \$62,000, Q3 = \$58,000, Q4 = \$71,000

\textbf{Solution:}
\begin{align*}
\text{Total Sales} &= Q1 + Q2 + Q3 + Q4 \\
&= 45 + 62 + 58 + 71 \\
&= 236 \text{ (in thousands)}
\end{align*}

\textbf{Answer:} Total sales for 2025 = \textbf{\$236,000}

\subsection*{Solution 2}
\textbf{Question:} By how much did sales increase from Q1 to Q4?

\textbf{Given:} Q1 = \$45,000, Q4 = \$71,000

\textbf{Solution:}
\begin{align*}
\text{Increase} &= Q4 - Q1 \\
&= 71 - 45 \\
&= 26 \text{ (in thousands)}
\end{align*}

\textbf{Answer:} Sales increased by \textbf{\$26,000}

\subsection*{Solution 3}
\textbf{Question:} Which employee worked the most hours?

\textbf{Given Table:}
\begin{center}
\begin{tabular}{|c|c|c|c|c|c|}
\hline
Employee & Alice & Bob & Carol & David & Emma \\
\hline
Hours & 40 & 38 & 42 & 35 & 45 \\
\hline
\end{tabular}
\end{center}

\textbf{Solution:} Comparing the hours: Alice (40), Bob (38), Carol (42), David (35), Emma (45).

The maximum value is 45 hours.

\textbf{Answer:} \textbf{Emma} worked the most hours with 45 hours.

\subsection*{Solution 4}
\textbf{Question:} How many total hours did all employees work?

\textbf{Solution:}
\begin{align*}
\text{Total Hours} &= 40 + 38 + 42 + 35 + 45 \\
&= 200 \text{ hours}
\end{align*}

\textbf{Answer:} Total hours = \textbf{200 hours}

\subsection*{Solution 5}
\textbf{Question:} On which day did the fitness center have the lowest number of sign-ups?

\textbf{Given Data:} Mon (8), Tue (12), Wed (15), Thu (11), Fri (18), Sat (22)

\textbf{Solution:} Comparing the values: The minimum is 8 sign-ups.

\textbf{Answer:} \textbf{Monday} had the lowest number of sign-ups with 8.

\subsection*{Solution 6}
\textbf{Question:} What was the total number of sign-ups for the week?

\textbf{Solution:}
\begin{align*}
\text{Total Sign-ups} &= 8 + 12 + 15 + 11 + 18 + 22 \\
&= 86 \text{ sign-ups}
\end{align*}

\textbf{Answer:} Total sign-ups = \textbf{86}

\subsection*{Solution 7}
\textbf{Question:} How many more pounds of apples were sold than carrots?

\textbf{Given:} Apples = 240 lbs, Carrots = 160 lbs

\textbf{Solution:}
\begin{align*}
\text{Difference} &= \text{Apples} - \text{Carrots} \\
&= 240 - 160 \\
&= 80 \text{ lbs}
\end{align*}

\textbf{Answer:} \textbf{80 more pounds} of apples were sold than carrots.

\subsection*{Solution 8}
\textbf{Question:} If broccoli costs \$2.50 per pound wholesale, what is the total cost of broccoli sold?

\textbf{Given:} Broccoli sold = 180 lbs, Cost = \$2.50/lb

\textbf{Solution:}
\begin{align*}
\text{Total Cost} &= \text{Pounds} \times \text{Cost per Pound} \\
&= 180 \times 2.50 \\
&= 450 \text{ dollars}
\end{align*}

\textbf{Answer:} Total cost = \textbf{\$450}

\subsection*{Solution 9}
\textbf{Question:} Which meal type had the highest satisfaction score?

\textbf{Given Scores:} Breakfast (8.2), Lunch (7.9), Dinner (8.6), Late Night (7.4)

\textbf{Solution:} The maximum score is 8.6.

\textbf{Answer:} \textbf{Dinner} had the highest satisfaction score at 8.6 out of 10.

\subsection*{Solution 10}
\textbf{Question:} What is the difference between the highest and lowest satisfaction scores?

\textbf{Given:} Highest = 8.6, Lowest = 7.4

\textbf{Solution:}
\begin{align*}
\text{Difference} &= 8.6 - 7.4 \\
&= 1.2 \text{ points}
\end{align*}

\textbf{Answer:} Difference = \textbf{1.2 points}

\newpage

\section{Part B: Box Plots and Data Interpretation (Solutions 11-20)}

\subsection*{Solution 11}
\textbf{Question:} What is the interquartile range (IQR)?

\textbf{Given:} Q1 = 62, Q3 = 85

\textbf{Solution:}
\begin{align*}
\text{IQR} &= Q3 - Q1 \\
&= 85 - 62 \\
&= 23
\end{align*}

\textbf{Answer:} IQR = \textbf{23}

\subsection*{Solution 12}
\textbf{Question:} How many of the data values fall between Q1 and Q3?

\textbf{Solution:} By definition, the interquartile range (the region between Q1 and Q3) contains the middle 50\% of the data values.

\textbf{Answer:} \textbf{50\%} (or the middle half) of the data values fall between Q1 and Q3.

\subsection*{Solution 13}
\textbf{Question:} What is the range of the data?

\textbf{Given:} Min = 12, Max = 78

\textbf{Solution:}
\begin{align*}
\text{Range} &= \text{Max} - \text{Min} \\
&= 78 - 12 \\
&= 66
\end{align*}

\textbf{Answer:} Range = \textbf{66}

\subsection*{Solution 14}
\textbf{Question:} What percentage of values are between the minimum and Q1?

\textbf{Solution:} By definition, Q1 (the first quartile) marks the point where 25\% of the data lies below it and 75\% lies above it. Therefore, values between the minimum and Q1 represent the lowest 25\% of the data.

\textbf{Answer:} \textbf{25\%} of values are between the minimum and Q1.

\subsection*{Solution 15}
\textbf{Question:} Identify any outliers (values more than 1.5 times the IQR below Q1 or above Q3).

\textbf{Given:} Min = 78\%, Q1 = 88\%, Median = 92\%, Q3 = 95\%, Max = 99\%

\textbf{Solution:}
First, calculate the IQR:
\begin{align*}
\text{IQR} &= Q3 - Q1 \\
&= 95 - 88 \\
&= 7\%
\end{align*}

Calculate outlier boundaries:
\begin{align*}
\text{Lower Boundary} &= Q1 - 1.5 \times \text{IQR} \\
&= 88 - 1.5(7) \\
&= 88 - 10.5 \\
&= 77.5\%
\end{align*}

\begin{align*}
\text{Upper Boundary} &= Q3 + 1.5 \times \text{IQR} \\
&= 95 + 1.5(7) \\
&= 95 + 10.5 \\
&= 105.5\%
\end{align*}

Check data: Min (78\%) < 77.5\%, so 78\% is a lower outlier. All other values fall within the boundaries.

\textbf{Answer:} \textbf{78\% is an outlier} (lower outlier).

\subsection*{Solution 16}
\textbf{Question:} How many quarters of the data lie above the median?

\textbf{Solution:} By definition, the median divides the data in half. The data above the median is represented by the range from the median to the maximum, which includes Q3 and the values above it. This represents 50\% or two quarters of the data. More specifically, the range from the median to Q3 is one quarter, and from Q3 to the maximum is another quarter.

\textbf{Answer:} \textbf{Two quarters} (or 50\%) of the data lie above the median.

\subsection*{Solution 17}
\textbf{Question:} What is the IQR for this data?

\textbf{Given:} Q1 = 2.5, Median = 4.0, Q3 = 5.5

\textbf{Solution:}
\begin{align*}
\text{IQR} &= Q3 - Q1 \\
&= 5.5 - 2.5 \\
&= 3.0 \text{ inches}
\end{align*}

\textbf{Answer:} IQR = \textbf{3.0 inches}

\subsection*{Solution 18}
\textbf{Question:} What is the lower outlier boundary?

\textbf{Given:} Q1 = 2.5, IQR = 3.0

\textbf{Solution:}
\begin{align*}
\text{Lower Outlier Boundary} &= Q1 - 1.5 \times \text{IQR} \\
&= 2.5 - 1.5(3.0) \\
&= 2.5 - 4.5 \\
&= -2.0 \text{ inches}
\end{align*}

\textbf{Answer:} Lower outlier boundary = \textbf{-2.0 inches}

\subsection*{Solution 19}
\textbf{Question:} What is Q1?

\textbf{Given:} IQR = 20, Q3 = 80

\textbf{Solution:}
\begin{align*}
\text{IQR} &= Q3 - Q1 \\
20 &= 80 - Q1 \\
Q1 &= 80 - 20 \\
Q1 &= 60
\end{align*}

\textbf{Answer:} Q1 = \textbf{60}

\subsection*{Solution 20}
\textbf{Question:} What is the upper outlier boundary?

\textbf{Given:} Q1 = 30, Q3 = 50, so IQR = 20

\textbf{Solution:}
\begin{align*}
\text{Upper Outlier Boundary} &= Q3 + 1.5 \times \text{IQR} \\
&= 50 + 1.5(20) \\
&= 50 + 30 \\
&= 80
\end{align*}

\textbf{Answer:} Upper outlier boundary = \textbf{80}

\newpage

\section{Part C: Combining Multiple Data Sources (Solutions 21-30)}

\subsection*{Solution 21}
\textbf{Question:} Which store had higher total sales over the three days?

\textbf{Given:}
\begin{center}
\begin{tabular}{|c|c|c|c|}
\hline
Store & Day 1 & Day 2 & Day 3 \\
\hline
Store A & \$500 & \$620 & \$580 \\
Store B & \$450 & \$510 & \$640 \\
\hline
\end{tabular}
\end{center}

\textbf{Solution:}
\begin{align*}
\text{Store A Total} &= 500 + 620 + 580 = 1700 \\
\text{Store B Total} &= 450 + 510 + 640 = 1600
\end{align*}

Store A: \$1,700 > Store B: \$1,600

\textbf{Answer:} \textbf{Store A} had higher total sales with \$1,700.

\subsection*{Solution 22}
\textbf{Question:} What was the combined total sales for both stores?

\textbf{Solution:}
\begin{align*}
\text{Combined Total} &= \text{Store A} + \text{Store B} \\
&= 1700 + 1600 \\
&= 3300 \text{ dollars}
\end{align*}

\textbf{Answer:} Combined total sales = \textbf{\$3,300}

\subsection*{Solution 23}
\textbf{Question:} What is the combined average score for both classes?

\textbf{Given:}
\begin{center}
\begin{tabular}{|c|c|c|}
\hline
Class & Number of Students & Average Score \\
\hline
Class A & 25 & 78 \\
Class B & 30 & 82 \\
\hline
\end{tabular}
\end{center}

\textbf{Solution:} Use the weighted average formula:
\begin{align*}
\text{Combined Average} &= \frac{\text{(Class A students} \times \text{Class A avg)} + \text{(Class B students} \times \text{Class B avg)}}{\text{Total students}} \\
&= \frac{(25 \times 78) + (30 \times 82)}{25 + 30} \\
&= \frac{1950 + 2460}{55} \\
&= \frac{4410}{55} \\
&= 80.18
\end{align*}

\textbf{Answer:} Combined average score = \textbf{80.18} (or approximately 80.2)

\subsection*{Solution 24}
\textbf{Question:} What is the overall average salary across all departments?

\textbf{Given:}
\begin{center}
\begin{tabular}{|c|c|c|}
\hline
Department & Employees & Avg. Salary \\
\hline
Sales & 20 & \$45,000 \\
Marketing & 15 & \$55,000 \\
Operations & 35 & \$50,000 \\
\hline
\end{tabular}
\end{center}

\textbf{Solution:} Use weighted average:
\begin{align*}
\text{Overall Avg. Salary} &= \frac{(20 \times 45000) + (15 \times 55000) + (35 \times 50000)}{20 + 15 + 35} \\
&= \frac{900000 + 825000 + 1750000}{70} \\
&= \frac{3475000}{70} \\
&= 49642.86
\end{align*}

\textbf{Answer:} Overall average salary = \textbf{\$49,642.86} (or approximately \$49,643)

\subsection*{Solution 25}
\textbf{Question:} Which gym had the greater increase from January to April?

\textbf{Given:}
\begin{center}
\begin{tabular}{|c|c|c|}
\hline
Month & Gym X & Gym Y \\
\hline
Jan & 150 & 200 \\
Feb & 165 & 210 \\
Mar & 180 & 205 \\
Apr & 195 & 220 \\
\hline
\end{tabular}
\end{center}

\textbf{Solution:}
\begin{align*}
\text{Gym X Increase} &= 195 - 150 = 45 \\
\text{Gym Y Increase} &= 220 - 200 = 20
\end{align*}

Gym X: 45 > Gym Y: 20

\textbf{Answer:} \textbf{Gym X} had the greater increase with 45 additional members.

\subsection*{Solution 26}
\textbf{Question:} What is the average monthly membership for Gym X over the four months?

\textbf{Solution:}
\begin{align*}
\text{Average Membership} &= \frac{150 + 165 + 180 + 195}{4} \\
&= \frac{690}{4} \\
&= 172.5
\end{align*}

\textbf{Answer:} Average monthly membership for Gym X = \textbf{172.5 members}

\subsection*{Solution 27}
\textbf{Question:} What is the total revenue for both regions across all quarters?

\textbf{Given:}
\begin{center}
\begin{tabular}{|c|c|c|c|c|}
\hline
Region & Q1 & Q2 & Q3 & Q4 \\
\hline
North & 12 & 15 & 18 & 22 \\
South & 10 & 14 & 16 & 20 \\
\hline
\end{tabular}
\end{center}

\textbf{Solution:}
\begin{align*}
\text{North Total} &= 12 + 15 + 18 + 22 = 67 \text{ million} \\
\text{South Total} &= 10 + 14 + 16 + 20 = 60 \text{ million} \\
\text{Combined Total} &= 67 + 60 = 127 \text{ million}
\end{align*}

\textbf{Answer:} Total revenue = \textbf{\$127 million}

\subsection*{Solution 28}
\textbf{Question:} Which quarter had the highest combined revenue from both regions?

\textbf{Solution:}
\begin{align*}
\text{Q1 Combined} &= 12 + 10 = 22 \\
\text{Q2 Combined} &= 15 + 14 = 29 \\
\text{Q3 Combined} &= 18 + 16 = 34 \\
\text{Q4 Combined} &= 22 + 20 = 42
\end{align*}

Q4 has the highest combined revenue at \$42 million.

\textbf{Answer:} \textbf{Q4} had the highest combined revenue at \$42 million.

\subsection*{Solution 29}
\textbf{Question:} What is the overall average hours of exercise for all respondents?

\textbf{Given:}
\begin{center}
\begin{tabular}{|c|c|c|}
\hline
Age Group & Sample Size & Avg. Hours \\
\hline
18-30 & 120 & 5.2 \\
31-45 & 100 & 4.1 \\
46+ & 80 & 2.8 \\
\hline
\end{tabular}
\end{center}

\textbf{Solution:} Use weighted average:
\begin{align*}
\text{Overall Avg. Hours} &= \frac{(120 \times 5.2) + (100 \times 4.1) + (80 \times 2.8)}{120 + 100 + 80} \\
&= \frac{624 + 410 + 224}{300} \\
&= \frac{1258}{300} \\
&= 4.19\overline{3}
\end{align*}

\textbf{Answer:} Overall average hours = \textbf{4.19 hours} (or approximately 4.2 hours)

\subsection*{Solution 30}
\textbf{Question:} How many total respondents were surveyed?

\textbf{Solution:}
\begin{align*}
\text{Total Respondents} &= 120 + 100 + 80 = 300
\end{align*}

\textbf{Answer:} Total respondents = \textbf{300}

\newpage

\section{Part D: Mean, Median, and Mode (Solutions 31-40)}

\subsection*{Solution 31}
\textbf{Question:} Find the mean of the dataset: 12, 15, 18, 14, 16

\textbf{Solution:}
\begin{align*}
\text{Mean} &= \frac{12 + 15 + 18 + 14 + 16}{5} \\
&= \frac{75}{5} \\
&= 15
\end{align*}

\textbf{Answer:} Mean = \textbf{15}

\subsection*{Solution 32}
\textbf{Question:} Find the median of the dataset: 24, 31, 19, 28, 25, 22, 29

\textbf{Solution:} First, arrange the data in order: 19, 22, 24, 25, 28, 29, 31

There are 7 values. The median is the middle value (position $\frac{7+1}{2} = 4$).

\textbf{Answer:} Median = \textbf{25}

\subsection*{Solution 33}
\textbf{Question:} Find the mode of the dataset: 5, 7, 5, 9, 5, 8, 7, 5

\textbf{Solution:} Count the frequency of each value:
\begin{itemize}
    \item 5 appears 4 times
    \item 7 appears 2 times
    \item 9 appears 1 time
    \item 8 appears 1 time
\end{itemize}

The mode is the value that appears most frequently.

\textbf{Answer:} Mode = \textbf{5}

\subsection*{Solution 34}
\textbf{Question:} A dataset has the values: 8, 12, 15, 20, 25. If a value of 100 is added, how does the median change?

\textbf{Solution:}
Original dataset (5 values): 8, 12, 15, 20, 25. Median = 15 (middle value)

New dataset (6 values): 8, 12, 15, 20, 25, 100. Ordered: 8, 12, 15, 20, 25, 100

For 6 values, the median is the average of the 3rd and 4th values:
\begin{align*}
\text{New Median} &= \frac{15 + 20}{2} = \frac{35}{2} = 17.5
\end{align*}

Change: 17.5 - 15 = 2.5 increase

\textbf{Answer:} The median increases from 15 to \textbf{17.5}, an increase of \textbf{2.5}.

\subsection*{Solution 35}
\textbf{Question:} The mean of five numbers is 24. If four of the numbers are 18, 22, 26, and 28, what is the fifth number?

\textbf{Solution:} If the mean of 5 numbers is 24, then:
\begin{align*}
\frac{18 + 22 + 26 + 28 + x}{5} &= 24 \\
\frac{94 + x}{5} &= 24 \\
94 + x &= 120 \\
x &= 26
\end{align*}

Verification: $\frac{18 + 22 + 26 + 28 + 26}{5} = \frac{120}{5} = 24$ 

\textbf{Answer:} The fifth number is \textbf{26}.

\subsection*{Solution 36}
\textbf{Question:} A dataset has six values with a median of 30. If the values in order are 15, 22, $x$, 35, 40, 48, what is $x$?

\textbf{Solution:} For 6 ordered values, the median is the average of the 3rd and 4th values:
\begin{align*}
\text{Median} &= \frac{x + 35}{2} \\
30 &= \frac{x + 35}{2} \\
60 &= x + 35 \\
x &= 25
\end{align*}

Verification: $\frac{25 + 35}{2} = \frac{60}{2} = 30$ 

\textbf{Answer:} $x = \textbf{25}$

\subsection*{Solution 37}
\textbf{Question:} Find the mean, median, and mode for: 42, 48, 42, 45, 50, 42, 48

\textbf{Solution:}

\textit{Mean:}
\begin{align*}
\text{Mean} &= \frac{42 + 48 + 42 + 45 + 50 + 42 + 48}{7} \\
&= \frac{317}{7} \\
&= 45.29 \text{ (approximately)}
\end{align*}

\textit{Median:} Ordered: 42, 42, 42, 45, 48, 48, 50
The middle value (4th position) is 45.

\textit{Mode:} Count frequencies:
\begin{itemize}
    \item 42 appears 3 times
    \item 48 appears 2 times
    \item 45 appears 1 time
    \item 50 appears 1 time
\end{itemize}

Mode = 42

\textbf{Answer:} Mean = \textbf{45.29}, Median = \textbf{45}, Mode = \textbf{42}

\subsection*{Solution 38}
\textbf{Question:} A teacher grades papers and records scores: 85, 92, 88, 90, 85, 92, 88, 95. What is the mode?

\textbf{Solution:} Count frequencies:
\begin{itemize}
    \item 85 appears 2 times
    \item 92 appears 2 times
    \item 88 appears 2 times
    \item 90 appears 1 time
    \item 95 appears 1 time
\end{itemize}

This dataset is trimodal (three modes): 85, 88, and 92 all appear equally (2 times each).

\textbf{Answer:} The modes are \textbf{85, 88, and 92} (trimodal dataset)

\subsection*{Solution 39}
\textbf{Question:} The mean age of a group of 8 people is 35 years. If a 9th person aged 50 joins the group, what is the new mean age?

\textbf{Solution:}
\begin{align*}
\text{Sum of 8 people's ages} &= 8 \times 35 = 280 \\
\text{Sum including 9th person} &= 280 + 50 = 330 \\
\text{New Mean} &= \frac{330}{9} \\
&= 36.\overline{6} \text{ (or approximately 36.67)}
\end{align*}

\textbf{Answer:} New mean age = \textbf{36.67 years} (or $\mathbf{36\frac{2}{3}}$ years)

\subsection*{Solution 40}
\textbf{Question:} A dataset has the values: 10, 15, 20, 25, 30, 35, 40. Which measure of central tendency (mean, median, or mode) best represents this data and why?

\textbf{Solution:}
\begin{align*}
\text{Mean} &= \frac{10 + 15 + 20 + 25 + 30 + 35 + 40}{7} = \frac{175}{7} = 25 \\
\text{Median} &= \text{Middle value} = 25 \\
\text{Mode} &= \text{None (no value repeats)}
\end{align*}

\textbf{Analysis:} The data is evenly spaced with no outliers and no repeated values. The mean and median are equal (both 25), and there is no mode. This is a symmetric distribution.

\textbf{Answer:} Either the \textbf{mean or median} is best, as they are equal and this symmetric, evenly-spaced distribution has no outliers or skewness. The mode cannot be used since no value repeats. The mean of 25 or median of 25 equally represent the center of this dataset.

\newpage

\section{Part E: Weighted Averages and Complex Problems (Solutions 41-50)}

\subsection*{Solution 41}
\textbf{Question:} A student's grade is calculated as follows: Homework (30\%), Midterm (40\%), Final (30\%). The student earned 85 on homework, 78 on the midterm, and 88 on the final. What is the weighted average grade?

\textbf{Solution:}
\begin{align*}
\text{Weighted Average} &= (0.30 \times 85) + (0.40 \times 78) + (0.30 \times 88) \\
&= 25.5 + 31.2 + 26.4 \\
&= 83.1
\end{align*}

\textbf{Answer:} Weighted average grade = \textbf{83.1}

\subsection*{Solution 42}
\textbf{Question:} A recipe calls for 2 cups of sugar (worth \$3 per cup) and 3 cups of flour (worth \$1 per cup). What is the average cost per cup of the final mixture?

\textbf{Solution:}
\begin{align*}
\text{Total cost of sugar} &= 2 \times 3 = 6 \text{ dollars} \\
\text{Total cost of flour} &= 3 \times 1 = 3 \text{ dollars} \\
\text{Total cost} &= 6 + 3 = 9 \text{ dollars} \\
\text{Total cups} &= 2 + 3 = 5 \text{ cups} \\
\text{Average cost per cup} &= \frac{9}{5} = 1.80 \text{ dollars}
\end{align*}

\textbf{Answer:} Average cost per cup = \textbf{\$1.80}

\subsection*{Solution 43}
\textbf{Question:} A portfolio contains 100 shares of Stock A (\$50/share) and 200 shares of Stock B (\$75/share). What is the weighted average price per share?

\textbf{Solution:}
\begin{align*}
\text{Value of Stock A} &= 100 \times 50 = 5000 \text{ dollars} \\
\text{Value of Stock B} &= 200 \times 75 = 15000 \text{ dollars} \\
\text{Total value} &= 5000 + 15000 = 20000 \text{ dollars} \\
\text{Total shares} &= 100 + 200 = 300 \\
\text{Weighted average price} &= \frac{20000}{300} = 66.\overline{6} \text{ dollars}
\end{align*}

\textbf{Answer:} Weighted average price per share = \textbf{\$66.67}

\subsection*{Solution 44}
\textbf{Question:} An employee worked 20 hours at \$15/hour and 10 hours at \$18/hour. What is the weighted average hourly wage?

\textbf{Solution:}
\begin{align*}
\text{Earnings at } \$15/\text{hr} &= 20 \times 15 = 300 \text{ dollars} \\
\text{Earnings at } \$18/\text{hr} &= 10 \times 18 = 180 \text{ dollars} \\
\text{Total earnings} &= 300 + 180 = 480 \text{ dollars} \\
\text{Total hours} &= 20 + 10 = 30 \text{ hours} \\
\text{Weighted average wage} &= \frac{480}{30} = 16 \text{ dollars/hour}
\end{align*}

\textbf{Answer:} Weighted average hourly wage = \textbf{\$16/hour}

\subsection*{Solution 45}
\textbf{Question:} A company's quarterly profits are: Q1: \$50,000 (weight 1), Q2: \$60,000 (weight 2), Q3: \$55,000 (weight 1.5), Q4: \$70,000 (weight 1.5). What is the weighted average quarterly profit?

\textbf{Solution:}
\begin{align*}
\text{Weighted Sum} &= (50000 \times 1) + (60000 \times 2) + (55000 \times 1.5) + (70000 \times 1.5) \\
&= 50000 + 120000 + 82500 + 105000 \\
&= 357500 \\
\text{Sum of weights} &= 1 + 2 + 1.5 + 1.5 = 6 \\
\text{Weighted average} &= \frac{357500}{6} \\
&= 59583.33
\end{align*}

\textbf{Answer:} Weighted average quarterly profit = \textbf{\$59,583.33}

\subsection*{Solution 46}
\textbf{Question:} A store sells three types of apples at different prices:
\begin{center}
\begin{tabular}{|c|c|c|}
\hline
Type & Pounds Sold & Price per Pound \\
\hline
Red & 80 & \$1.50 \\
Green & 60 & \$1.75 \\
Golden & 40 & \$2.00 \\
\hline
\end{tabular}
\end{center}
What is the weighted average price per pound of apples sold?

\textbf{Solution:}
\begin{align*}
\text{Revenue from Red} &= 80 \times 1.50 = 120 \text{ dollars} \\
\text{Revenue from Green} &= 60 \times 1.75 = 105 \text{ dollars} \\
\text{Revenue from Golden} &= 40 \times 2.00 = 80 \text{ dollars} \\
\text{Total revenue} &= 120 + 105 + 80 = 305 \text{ dollars} \\
\text{Total pounds} &= 80 + 60 + 40 = 180 \text{ pounds} \\
\text{Weighted average price} &= \frac{305}{180} = 1.694\text{4} \text{ dollars/pound}
\end{align*}

\textbf{Answer:} Weighted average price per pound = \textbf{\$1.69} (or \$1.694)

\subsection*{Solution 47}
\textbf{Question:} A class with 25 students has an average test score of 82. Another class with 20 students has an average test score of 88. What is the combined average for both classes?

\textbf{Solution:}
\begin{align*}
\text{Total for Class 1} &= 25 \times 82 = 2050 \\
\text{Total for Class 2} &= 20 \times 88 = 1760 \\
\text{Combined total} &= 2050 + 1760 = 3810 \\
\text{Combined students} &= 25 + 20 = 45 \\
\text{Combined average} &= \frac{3810}{45} = 84.67
\end{align*}

\textbf{Answer:} Combined average = \textbf{84.67}

\subsection*{Solution 48}
\textbf{Question:} A weighted average calculation uses weights of 0.2, 0.3, 0.25, and 0.25 for values 60, 75, 80, and 90 respectively. What is the weighted average?

\textbf{Solution:}
\begin{align*}
\text{Weighted Average} &= (0.2 \times 60) + (0.3 \times 75) + (0.25 \times 80) + (0.25 \times 90) \\
&= 12 + 22.5 + 20 + 22.5 \\
&= 77
\end{align*}

Note: Verify that weights sum to 1: $0.2 + 0.3 + 0.25 + 0.25 = 1.0$ 

\textbf{Answer:} Weighted average = \textbf{77}

\subsection*{Solution 49}
\textbf{Question:} In a survey, 40\% of respondents chose Option A (satisfaction rating 8), 35\% chose Option B (satisfaction rating 6), and 25\% chose Option C (satisfaction rating 5). What is the weighted average satisfaction rating?

\textbf{Solution:}
\begin{align*}
\text{Weighted Average} &= (0.40 \times 8) + (0.35 \times 6) + (0.25 \times 5) \\
&= 3.2 + 2.1 + 1.25 \\
&= 6.55
\end{align*}

Note: Verify that percentages sum to 100\%: $40 + 35 + 25 = 100$ 

\textbf{Answer:} Weighted average satisfaction rating = \textbf{6.55}

\subsection*{Solution 50}
\textbf{Question:} A company manufactures a product in two facilities:
\begin{center}
\begin{tabular}{|c|c|c|}
\hline
Facility & Units Produced & Cost per Unit \\
\hline
Facility A & 500 & \$8 \\
Facility B & 700 & \$6 \\
\hline
\end{tabular}
\end{center}
What is the weighted average cost per unit?

\textbf{Solution:}
\begin{align*}
\text{Total cost at Facility A} &= 500 \times 8 = 4000 \text{ dollars} \\
\text{Total cost at Facility B} &= 700 \times 6 = 4200 \text{ dollars} \\
\text{Combined total cost} &= 4000 + 4200 = 8200 \text{ dollars} \\
\text{Total units produced} &= 500 + 700 = 1200 \text{ units} \\
\text{Weighted average cost} &= \frac{8200}{1200} = 6.8\overline{3} \text{ dollars/unit}
\end{align*}

\textbf{Answer:} Weighted average cost per unit = \textbf{\$6.83}

\newpage

\section{Summary}

This solutions document provides complete step-by-step worked solutions for all 50 questions covering:

\begin{itemize}
    \item \textbf{Part A (Questions 1-10):} Reading and interpreting tables — basic data extraction and calculations
    \item \textbf{Part B (Questions 11-20):} Box plots and data interpretation — quartiles, IQR, and outlier detection
    \item \textbf{Part C (Questions 21-30):} Combining multiple data sources — weighted averages and multi-group analysis
    \item \textbf{Part D (Questions 31-40):} Mean, median, mode — measures of central tendency and their applications
    \item \textbf{Part E (Questions 41-50):} Weighted averages and complex problems — real-world applications
\end{itemize}

All calculations have been verified, and the solutions follow the mathematical principles covered in the GED Math lessons on graphs, tables, and averages.

\end{document}